\subsection{Funcionalidades Clave e Innovación Central}

La plataforma Mi Trufi se sustenta en cuatro pilares funcionales diseñados para responder con precisión a las fallas operacionales identificadas en el diagnóstico.

\subsubsection{La Función Insignia: <<Levantar la Mano Virtual>>}

La característica distintiva del proyecto radica en digitalizar la señal física de parada mediante la funcionalidad de \emph{levantar la mano virtual}. El mecanismo opera de la siguiente manera:

\begin{enumerate}[label=\arabic*.]
    \item El pasajero abre la aplicación móvil, selecciona la línea o trayecto y pulsa el control para marcar su parada. La aplicación captura las coordenadas precisas del usuario mediante el receptor GPS del teléfono o le permite ajustar manualmente el marcador en el mapa.
    \item El servidor procesa la posición del pasajero y calcula qué micro en servicio se aproxima hacia ese punto dentro del corredor vial.
    \item La aplicación del chofer proyecta en su mapa un indicador visual de parada solicitada junto con una señal auditiva discreta, informando la distancia al punto de recojo.
    \item Si el chofer dispone de espacio, el abordaje se concreta de forma coordinada, ahorrando maniobras y reduciendo la ansiedad de ambas partes.
\end{enumerate}

\subsubsection{Telemetría y Monitoreo Satelital (GPS) en Tiempo Real}

A diferencia de los sistemas de transporte masivo que requieren costosos equipos embarcados dedicados, Mi Trufi aprovecha el hardware nativo de los teléfonos inteligentes de los conductores:

\begin{itemize}
    \item \textbf{Transmisión adaptativa de coordenadas:} El módulo del chofer reporta paquetes livianos de geolocalización a intervalos periódicos optimizados para no saturar el plan de datos móviles ni drenar prematuramente la batería del teléfono.
    \item \textbf{Cálculo de tiempos estimados de llegada (ETA):} La plataforma proyecta en los dispositivos de los pasajeros la distancia y el tiempo aproximado en que el micro llegará a la parada elegida, considerando la velocidad media de avance.
\end{itemize}

\subsubsection{Indicador Dinámico de Capacidad y Ocupación}

Para resolver el problema del sobrecupo en paradas intermedias, el chofer dispone de un selector táctil simple que no interfiere con la conducción:

\begin{itemize}
    \item \textbf{Estado Disponible:} Indica que la unidad cuenta con amplio espacio para recoger pasajeros a lo largo de la ruta.
    \item \textbf{Estado Cupo Limitado ($N$ Asientos):} Permite al chofer definir cuántos asientos libres quedan (por ejemplo, 1, 2 o 3 asientos), de modo que la aplicación de los pasajeros solo admita ese número exacto de solicitudes de parada.
    \item \textbf{Estado Lleno:} Alerta a los pasajeros en el mapa de que la unidad circula a capacidad máxima, permitiéndoles orientar su expectativa hacia el siguiente turno de salida de 15 minutos.
\end{itemize}

\subsubsection{Control y Despacho Automatizado de Turnos de 15 Minutos}

El panel web de secretaría sustituye la libreta de anotaciones manuales por un sistema de colas digitales:

\begin{itemize}
    \item Programa automáticamente la grilla horaria de salidas con intervalos regulares de 15 minutos.
    \item Registra la hora exacta en que cada chofer inicia su trayecto y valida que el vehículo en circulación corresponda al turno asignado.
    \item Detecta retrasos anormales en el trayecto de vuelta o salidas tardías, generando métricas de cumplimiento para la directiva de la línea.
\end{itemize}

\subsubsection{Diagramas de Interacción y Operatividad del Sistema}

La Figura~\ref{fig:flujo_actores} expone la articulación entre los perfiles de usuario y el servidor central, destacando la separación entre la plataforma web de secretaría y las aplicaciones móviles de campo.

\begin{figure}[htbp]
\centering
\resizebox{0.9\textwidth}{!}{%
\begin{tikzpicture}[
    actorBox/.style={rectangle, draw=black, fill=white, minimum width=3.4cm, minimum height=1.3cm, align=center, font=\small\bfseries, thick},
    serverBox/.style={rectangle, draw=black, fill=gray!10, minimum width=4.5cm, minimum height=2.2cm, align=center, font=\small\bfseries, thick},
    arrowLine/.style={-{Stealth[length=3mm]}, thick, draw=black}
]

    % Nodo Central
    \node[serverBox] (servidor) at (0, 0) {
        Servidor Central Mi Trufi\\
        \vspace{0.15cm}
        \normalfont\footnotesize 
        (Gestión de Rutas, Turnos,\\
        Telemetría y Notificaciones)
    };

    % Nodo Superior: Secretaría
    \node[actorBox, above=2.2cm of servidor] (secretaria) {
        Secretaría / Control\\
        \vspace{0.08cm}
        \normalfont\footnotesize (Plataforma Web)\\
        \normalfont\footnotesize Turnos de 15 min y Despacho
    };

    % Nodo Inferior Izquierdo: Chofer
    \node[actorBox, below left=2.2cm and 1.2cm of servidor] (chofer) {
        Chofer (Conductor)\\
        \vspace{0.08cm}
        \normalfont\footnotesize (Aplicación Móvil)\\
        \normalfont\footnotesize GPS y Estado de Cupo
    };

    % Nodo Inferior Derecho: Pasajero
    \node[actorBox, below right=2.2cm and 1.2cm of servidor] (pasajero) {
        Pasajero (Usuario)\\
        \vspace{0.08cm}
        \normalfont\footnotesize (Aplicación Móvil)\\
        \normalfont\footnotesize Levantado de Mano y Mapa
    };

    % Flujos de información
    \draw[arrowLine] (secretaria.south) -- (servidor.north) 
        node[midway, right, font=\footnotesize] {Registro de salidas y turnos};

    \draw[arrowLine] (chofer.north) -- (servidor.south west) 
        node[midway, left=0.1cm, font=\footnotesize, align=right] {Ubicación GPS y\\capacidad del micro};

    \draw[arrowLine] (pasajero.north) -- (servidor.south east) 
        node[midway, right=0.1cm, font=\footnotesize, align=left] {Solicitud de parada\\(Mano levantada)};

    \draw[arrowLine, dashed] (servidor.south west) -- (chofer.east) 
        node[midway, below=0.1cm, font=\scriptsize] {Avisos de pasajeros};

    \draw[arrowLine, dashed] (servidor.south east) -- (pasajero.west) 
        node[midway, below=0.1cm, font=\scriptsize] {Ubicación de unidades};

\end{tikzpicture}%
}
\caption{Flujo general de información e interacción entre los actores del sistema Mi Trufi.}
\label{fig:flujo_actores}
\end{figure}

Asimismo, en la Figura~\ref{fig:proceso_mano} se ilustra la secuencia operativa temporal desde el momento en que un usuario consulta el servicio hasta la concreción del abordaje.

\begin{figure}[htbp]
\centering
\resizebox{\textwidth}{!}{%
\begin{tikzpicture}[
    stepNode/.style={rectangle, draw=black, fill=white, minimum width=3.3cm, minimum height=1.6cm, align=center, font=\footnotesize, thick},
    arrowNode/.style={-{Stealth[length=3mm]}, thick, draw=black}
]

    \node[stepNode] (p1) at (0,0) {
        \textbf{Paso 1: Consulta}\\
        El pasajero visualiza el\\mapa y la proximidad\\de las unidades en ruta.
    };

    \node[stepNode, right=1.2cm of p1] (p2) {
        \textbf{Paso 2: Solicitud}\\
        El pasajero activa\\<<Levantar la mano>>\\marcando su ubicación.
    };

    \node[stepNode, right=1.2cm of p2] (p3) {
        \textbf{Paso 3: Notificación}\\
        El chofer recibe el punto\\en pantalla y confirma si\\dispone de asientos libres.
    };

    \node[stepNode, below=1.4cm of p2] (p4) {
        \textbf{Paso 4: Abordaje Coordinado}\\
        La unidad recoge al usuario de forma previsible, optimizando tiempo y detenciones.
    };

    \draw[arrowNode] (p1) -- (p2);
    \draw[arrowNode] (p2) -- (p3);
    \draw[arrowNode] (p3.south) |- (p4.east);

\end{tikzpicture}%
}
\caption{Secuencia operativa de la función para solicitar parada mediante la plataforma.}
\label{fig:proceso_mano}
\end{figure}
